%% ma: L10-B01-0034
%% lop: 10
%% chuong: 1
%% bai: 1
%% dang: 10.1.1
%% loai: TLN
%% muc: TH
%% dap_an: "2"
%% nguon: "Ảnh giáo viên gửi 10/10/2026 (ThS. Nguyễn Hoàng Việt, Chương 1 Mệnh đề), B-Đề 2, Phần 3 Câu 18"
%% kien_thuc: "Bài 1 (mệnh đề, mệnh đề chứa biến)"
%% trang_thai: cho-duyet
%% ly_do: "Dạng toán mới đề xuất 10.1.1: nhận biết câu là mệnh đề, câu không phải mệnh đề"
%% ghi_chu: "Đề gốc không có đáp án; đáp số 2 do tự giải theo định nghĩa mệnh đề toán học"
\begin{ex}
Trong các mệnh đề sau, có bao nhiêu mệnh đề toán học?
\begin{enumerate}[label=\alph*)]
\item $A\colon$ ``Tích hai số thực trái dấu là một số thực âm''.
\item $B\colon$ ``Mọi số tự nhiên đều là số dương''.
\item $C\colon$ ``Có sự sống ngoài Trái Đất''.
\item $D\colon$ ``Ngày $1$ tháng $5$ là ngày Quốc tế Lao động''.
\end{enumerate}
\shortans{2}
\loigiai{Mệnh đề toán học là mệnh đề có nội dung toán học. $A$ (đúng) và $B$ (sai vì số $0$ không dương) là mệnh đề toán học. $C$ chưa xác định được đúng sai; $D$ là mệnh đề nhưng không có nội dung toán học. Có $2$ mệnh đề toán học.}
\end{ex}
